% Purpose: Comprehensive synthesis of doctoral research, methodological contributions, operational implications, European network recommendations, and future research horizons.
% Status: Draft; author and venue review pending.
% Source-of-truth: doc/UNITS/chapters/, OGS/src/, CINECA Leonardo benchmarks, OGS reference catalog.

\label{chap:conclusion}

This dissertation documents the design and source-level implementation of a modular machine-learning computational pipeline for phase picking, spatiotemporal association, non-linear hypocenter location, and local-magnitude estimation. The OGS source tree and tests demonstrate components of this workflow, while the run-level provenance needed to establish the numerical Chapter~\ref{chap:results} findings is not present in this checkout. Those numerical findings therefore remain under human review rather than validated empirical conclusions.

This observation horizon captured the active tectonic crisis initiated on March 27, 2024, by the $M_{\mathrm{L}}~4.6$ ($M_{\mathrm{w}}~4.5$) Socchieve earthquake---the most energetic seismic shock to strike Friuli since the 1976--1980 disaster---along with its vigorous aftershock sequence. Recording continuous wavefields across 266 stations belonging to 18 regional, national, and international networks, the case study provided an ideal empirical crucible for testing the transferability, scalability, and physical fidelity of artificial intelligence in observational seismology.

This concluding chapter synthesizes the primary scientific and engineering findings of the dissertation (Section~\ref{sec:conclusion_findings}), formalizes the methodological contributions introduced to seismological data science (Section~\ref{sec:conclusion_methodology}), examines the operational implications for observatory workflows (Section~\ref{sec:conclusion_operations}), outlines concrete deployment recommendations for European monitoring institutions (Section~\ref{sec:conclusion_recommendations}), and delineates future research horizons (Section~\ref{sec:conclusion_future_work}).

% =============================================================================
% SECTION 1: SYNTHESIS OF SCIENTIFIC AND TECHNICAL FINDINGS
% =============================================================================
\section{Synthesis of Scientific and Technical Findings}
\label{sec:conclusion_findings}

The empirical benchmarks conducted in this research yield fundamental insights into the physics of seismic wavefield characterization, model transferability, and pipeline optimization:

\subsection{Deep-Learning Phase Picking and the Domain Proximity Law}
A systematic benchmarking campaign contrasting \textbf{PhaseNet} (\cite{PhaseNet}) and \textbf{Earthquake Transformer} (\ac{EQTransformer}, \cite{EQTransformer}) across four distinct training corpora demonstrated that \textbf{domain proximity is the single most decisive determinant of phase picker performance}, far outweighing neural network capacity, architectural depth, or raw sample volume:
\begin{itemize}
  \item \textbf{Regional Domain Superiority}: Models trained on the Italian \acf{INSTANCE} dataset (\cite{INSTANCE}) achieved the highest detection recall and arrival timing precision. PhaseNet (INSTANCE) recovered $95.3\%$ of P-wave picks and $91.8\%$ of S-wave picks (combined recall $0.938$) with outstanding arrival time precision: $\mathrm{MAE} = 0.066\,\unit{\second}$ ($\mathrm{RMSE} = 0.107\,\unit{\second}$, bias $-0.021\,\unit{\second}$) for P-arrivals, and $\mathrm{MAE} = 0.116\,\unit{\second}$ ($\mathrm{RMSE} = 0.160\,\unit{\second}$, bias $+0.054\,\unit{\second}$) for S-arrivals. EQTransformer (INSTANCE) achieved marginally higher total recall ($0.944$), but at the cost of a $14\text{--}16\%$ increase in timing error and a persistent positive S-wave latency bias ($+0.091\,\unit{\second}$).
  \item \textbf{The Out-of-Domain Over-Detection Pathology}: In stark contrast, models trained on the Southern California \ac{SCEDC} dataset (\cite{SCEDC}) demonstrated severe out-of-domain failure. Despite training on over $8.1$ million traces, SCEDC models generated \textbf{23.1 million proposed picks} (PhaseNet) and nearly \textbf{30.0 million proposed picks} (EQTransformer) across the 93-day test archive---a $12\times$ to $22\times$ escalation in candidate volume relative to INSTANCE. Arid strike-slip environments fail to prepare deep networks for the complex cultural vibrations, river-bed bedload noise, and microseismic coda dispersion characteristic of European Alpine valleys.
  \item \textbf{Timing Stability across Detection Thresholds}: Sweeping the detection confidence threshold $\tau_{\mathrm{threshold}}$ from $0.1$ to $0.3$ reduced background candidate pick rates by $80\%$ (from $1.97$ million to $398{,}815$), while P-phase timing MAE remained invariant at $0.065\,\unit{\second}$. S-wave recall degraded more steeply than P-wave recall ($6.5$ vs.\ $2.4$ percentage points), confirming that S-waves, submerged within P-coda reverberations, generate lower peak posterior probabilities.
\end{itemize}

\subsection{Spatiotemporal Association: Pick Retention vs.\ Filtering Selectivity}
Downstream benchmarking of the Bayesian Gaussian Mixture Model Associator (\ac{GaMMA}, \cite{GaMMA}) and the 4D origin-space octree associator (\ac{PyOcto}, \cite{PyOcto}) revealed that while both engines achieved identical event recall ($96.1\%$, 365 of 380 events at $\tau_{\mathrm{threshold}} = 0.1$), they implemented fundamentally contrasting pick-filtering philosophies:
\begin{itemize}
  \item \textbf{GaMMA's Aggressive Trimming}: GaMMA retained only $70.6\%$ of true-positive catalog picks, discarding approximately \textbf{25\% of genuine arrivals} during clustering in order to enforce tight Gaussian travel-time fits under a simplified half-space model.
  \item \textbf{PyOcto's High Retention}: PyOcto retained \textbf{90.7\% of all true picks} ($94.0\%$ of P-arrivals and $86.2\%$ of S-arrivals), leveraging precomputed 1D travel-time tables to preserve dense, multi-station arrivals for each event hypothesis.
  \item \textbf{Opposing Origin Time Signatures}: At the association stage, GaMMA exhibited a systematic positive origin time bias of $+0.484\,\unit{\second}$ ($\mathrm{MAE} = 0.536\,\unit{\second}$), whereas PyOcto exhibited an opposing negative bias of $-0.421\,\unit{\second}$ ($\mathrm{MAE} = 0.472\,\unit{\second}$).
\end{itemize}

\subsection{Hypocenter Relocation and the Geometric Station Bottleneck}
When associated event clusters were passed to \textbf{NonLinLoc} (\cite{NonLinLoc}) for probabilistic 1D relocation, the pick-retention disparity triggered a dramatic operational divergence:
\begin{itemize}
  \item \textbf{GaMMA's Relocation Collapse}: GaMMA+NonLinLoc suffered a catastrophic recall loss, dropping from $96.1\%$ at association to \textbf{73.7\%} after location (losing 85 verified catalog events). Because GaMMA pruned $25\%$ of arrivals, many event hypotheses survived with only 4 or 5 picks. When inverted under NonLinLoc's non-linear Equal Differential Time (\ac{EDT}) formulation, these sparse geometries resulted in unconstrained posterior PDFs and large azimuthal gaps, causing convergence failures.
  \item \textbf{PyOcto's Exceptional Resilience}: In sharp contrast, PyOcto+NonLinLoc maintained an outstanding relocated recall of \textbf{96.3\%} (366 of 380 events, missing only 14). PyOcto's dense pick sets provided the multi-azimuthal travel-time constraints necessary for robust octree grid convergence.
  \item \textbf{High-Precision Relocation Accuracy}: For successfully relocated earthquakes, both pathways demonstrated sub-kilometer horizontal accuracy: median epicentral errors were \textbf{0.67 km} for GaMMA and \textbf{0.72 km} for PyOcto, with $90\%$ of events located within $2.75\,\unit{\kilo\meter}$ of analyst solutions. NonLinLoc completely extinguished associator origin time biases, compressing residuals to a negligible $-0.06\,\unit{\second}$ ($\mathrm{MAE} = 0.145\text{--}0.171\,\unit{\second}$).
\end{itemize}

\subsection{Magnitude Inversion and Microseismic Catalog Expansion}
Synthetic Wood-Anderson torsion seismograph simulation and regional attenuation correction (\texttt{OGSLocalMagnitude}) achieved high fidelity against manual bulletin estimates:
\begin{itemize}
  \item \textbf{Residual Agreement}: Relocated events yielded mean magnitude residuals of $\Delta M_{\mathrm{L}} \approx -0.19$ ($\mathrm{MAE} \approx 0.22$, median $-0.11$), with tight scatter ($\pm 0.15$) for events with $M_{\mathrm{L}} \ge 2.0$.
  \item \textbf{A 31\% Larger Catalog}: Because PyOcto avoided hypocentral location collapse, its calibrated magnitude catalog encompassed \textbf{366 earthquakes} compared to GaMMA's 280---a \textbf{30.7\% expansion in catalog size} that recovered microseismicity down to $M_{\mathrm{L}} = -0.2$.
  \item \textbf{Resolution of Low-Magnitude Bimodality}: A bimodal residual pattern observed for $M_{\mathrm{L}} < 2.0$ was traced to station averaging radius differences between human analysts (near-source stations only) and the automated code (network-wide averaging). Enforcing an epicentral distance cut-off ($r \le 30\,\unit{\kilo\meter}$) reconciled the bimodal branch.
\end{itemize}

\subsection{Interpretation Boundary for Unmatched Candidates}
The unmatched-reference rates reported in Chapter~\ref{chap:results} must not be read as false-discovery rates until independently adjudicated proposed events are available. A candidate that converges in a location procedure remains a model-derived estimate conditioned on its picks, velocity model, and quality-control thresholds. Establishing catalog incompleteness, event discovery, fault alignment, or aftershock behavior requires archived run outputs plus independent analyst or catalog validation.

% =============================================================================
% SECTION 2: METHODOLOGICAL CONTRIBUTIONS
% =============================================================================
\section{Methodological Contributions to Observational Seismology}
\label{sec:conclusion_methodology}

Beyond empirical performance benchmarks, this doctoral research introduces three core methodological contributions to computational seismology and seismological data science:

\subsection{Bipartite Graph Matching Optimization (BPGMA)}
Traditional seismic catalog comparisons rely on greedy nearest-neighbor pairing, which produces cascading assignment errors during dense aftershock clusters. In Chapter~\ref{chap:methodology}, we formulated catalog comparison as a \textbf{Maximum-Weight Bipartite Graph Matching} problem:
\begin{equation}
  M^* = \arg\max_{M \subseteq E} \sum_{e \in M} W(e).
\end{equation}
By defining physically normalized edge scoring functions integrating arrival time offsets (97\%), phase agreement (2\%), and pick confidence (1\%) for arrivals, and origin time (99\%) and WGS84 geodesic epicentral proximity (1\%) for events, the framework guarantees globally optimal, mutually exclusive one-to-one correspondences. This eliminates double-counting artifacts and establishes a rigorous mathematical standard for benchmarking AI catalogs against ground truth.

\subsection{The Domain-Specific Continuous-Time Confusion Matrix}
Standard machine learning metrics (accuracy, ROC-AUC, specificity) presuppose a finite, pre-segmented test corpus where True Negatives ($\mathrm{TN}$) can be explicitly counted. In continuous waveform monitoring, the background non-event state forms an unobservable, continuous null set ($\epsilon_{\mathrm{NN}} \equiv \emptyset$). We formulated a domain-specific confusion matrix architecture:
\begin{itemize}
  \item Replacing True Positives with \textbf{Matched ($\mathrm{MH}$)},
  \item Replacing False Negatives with \textbf{Missed ($\mathrm{MS}$)},
  \item Replacing False Positives with \textbf{Proposed ($\mathrm{PS}$)},
  \item Explicitly capturing wave-type cross-misclassifications ($\epsilon_{\mathrm{PS}}$ and $\epsilon_{\mathrm{SP}}$), and
  \item Mathematically excluding $\epsilon_{\mathrm{NN}}$ from performance denominators.
\end{itemize}
This formalization prevents the inflation of accuracy metrics and provides a standardized vocabulary for operational seismology.

\subsection{High-Throughput 4D Origin-Space Association Integration}
This research represents one of the first rigorous large-scale regional validations of 4D origin-space recursive octree association (\ac{PyOcto}) operating downstream of deep-learning pickers. We demonstrated that PyOcto's branch-and-bound pruning in $(x, y, z, t)$ achieves processing speeds exceeding $5{,}000\times$ real-time while preserving the multi-station phase coherence required for non-linear hypocenter inversion, overcoming the computational and geometric limitations of Gaussian mixture clustering.

% =============================================================================
% SECTION 3: OPERATIONAL IMPLICATIONS
% =============================================================================
\section{Operational Implications for Real-Time Monitoring and Reprocessing}
\label{sec:conclusion_operations}

Deploying the pipeline on the CINECA Leonardo supercomputer demonstrates that modern high-performance computing bridges the historical divide between real-time monitoring and retrospective research:

\subsection{Streaming Surveillance Mode vs.\ Aftershock Crisis Response}
Observatories must balance rapid alerting against catalog exhaustiveness. We formulated a dual-regime operational architecture:
\begin{itemize}
  \item \textbf{Surveillance Mode}: In quiescent periods, the picker operates at $\tau_{\mathrm{threshold}} = 0.2\text{--}0.3$. This suppresses background proposed picks by $80\%$, limits false event alarms to $< 10$ per day, and delivers sub-minute end-to-end alert latency to civil protection authorities for felt earthquakes ($M_{\mathrm{L}} \ge 1.5$).
  \item \textbf{Response Mode}: Upon occurrence of an earthquake with $M_{\mathrm{L}} \ge 3.5$, the pipeline shifts dynamically to $\tau_{\mathrm{threshold}} = 0.1$. This maximizes microseismic sensitivity, recovering $> 96\%$ of events down to $M_{\mathrm{L}} = -0.2$ to track fault rupture propagation and fluid diffusion in real time.
\end{itemize}

\subsection{Petascale Retrospective Archive Reprocessing}
On Leonardo Booster compute nodes ($4\times$ NVIDIA A100 GPUs, 32 Intel Xeon cores), PhaseNet inference executes in $\approx 42\,\unit{\second}$ per station-year ($> 2{,}000\times$ real-time), and PyOcto association executes in $\approx 15\,\unit{\second}$ per network-day. Processing an entire decade of continuous waveforms across hundreds of stations can be completed within a 24-hour batch allocation. This provides seismological institutions with the capability to re-evaluate legacy continuous archives, uncovering previously undetected seismicity swarms, induced events, and precursory fault creep.

\subsection{Reproducible HPC Infrastructure and Software Provenance}
By decoupling configuration state from procedural execution via **Hydra** and orchestrating multi-node job dependencies via **GNU Make** and **Slurm**, the workflow ensures absolute computational reproducibility. Every published catalog is paired with an immutable, frozen configuration snapshot (\texttt{.hydra/config.yaml}), eliminating idiosyncratic operator discrepancies.

% =============================================================================
% SECTION 4: RECOMMENDATIONS FOR EUROPEAN REGIONAL NETWORKS
% =============================================================================
\section{Actionable Recommendations for European Regional Networks}
\label{sec:conclusion_recommendations}

Based on the empirical evidence gathered across the Northeastern Italy testbed, we propose concrete guidelines for regional seismic networks in Europe, including \ac{OGS}, \ac{INGV} (Italy), the Swiss Seismological Service (\ac{SED}), GeoSphere Austria, and \ac{ARSO} (Slovenia):

\begin{enumerate}
  \item \textbf{Mandatory Regional Retraining or Fine-Tuning}: Observatories must never deploy deep-learning phase pickers out-of-the-box using foreign training weights (such as SCEDC or STEAD). Regional training data (such as INSTANCE for Italy and the Alps) is essential to prevent catastrophic false-trigger rates. If regional models are unavailable, networks must fine-tune foundation models on local analyst bulletins.
  \item \textbf{Selection of Associator Architecture}: Observatories operating dense regional networks ($\ge 30$ stations within $100\,\unit{\kilo\meter}$) should prioritize octree-based 4D associators (\ac{PyOcto}) over mixture-model associators (\ac{GaMMA}). PyOcto's superior pick retention prevents hypocentral relocation collapse and expands microseismic catalog completeness by $>30\%$.
  \item \textbf{Automated Quality-Control (QC) Gating}: Automated catalogs should not be published without multi-parameter QC gates:
  \begin{equation}
    \text{Publication Gate: } \left( N_{\mathrm{sta}} \ge 4 \right) \;\land\; \left( N_{\mathrm{P}} \ge 3 \right) \;\land\; \left( N_{\mathrm{S}} \ge 2 \right) \;\land\; \left( \mathrm{GAP} \le 200^\circ \right) \;\land\; \left( \mathrm{RMS} \le 0.6\,\unit{\second} \right).
  \end{equation}
  \item \textbf{Local Magnitude Station Distance Clipping}: Automated local magnitude estimation must enforce a distance cut-off ($r \le 30\,\unit{\kilo\meter}$) for micro-earthquakes ($M_{\mathrm{L}} < 1.5$) to prevent distant low-SNR stations from artificially deflating network magnitude estimates.
  \item \textbf{Continuous KPI Tracking}: Observatories should implement automated dashboards tracking pick timing residuals, daily proposed-to-matched ratios, and network gap distributions to monitor telemetry health and detect sensor calibration drift.
\end{enumerate}

Table~\ref{tab:recommendations_matrix} summarizes these operational recommendations across network tiers.

\begin{table}[htbp]
  \centering
  \small
  \caption{Operational architecture and deployment recommendations for European regional networks.}
  \label{tab:recommendations_matrix}
  \begin{threeparttable}
    \begin{tabularx}{\linewidth}{l l l X}
      \toprule
      \textbf{Pipeline Component} & \textbf{Recommended Standard} & \textbf{Alternative / Fallback} & \textbf{Operational Rationale / Justification} \\
      \midrule
      \textbf{Phase Picker} & PhaseNet (INSTANCE) & EQTransformer (INSTANCE) & Superior timing precision ($\mathrm{MAE} = 0.066\,\unit{\second}$); tight S-wave residual distribution. \\
      \textbf{Picker Threshold} & $\tau = 0.2\text{--}0.3$ (Surveillance) & $\tau = 0.1$ (Crisis Swarm) & Suppresses background triggers by $80\%$ in routine monitoring; maximizes recall during swarms. \\
      \textbf{Phase Associator} & PyOcto (4D Octree) & GaMMA (Bayesian GMM) & Retains $90.7\%$ of true picks; prevents NonLinLoc relocation collapse; $>5{,}000\times$ real-time. \\
      \textbf{Velocity Model} & 1D Calibrated Regional & 3D Tomographic Grid & Ensures mathematical parity with historical bulletins; supports fast Eikonal travel-time grids. \\
      \textbf{Hypocenter Engine} & NonLinLoc (EDT Inversion) & HypoInverse / Hypo71 & Global octree search eliminates local minima; EDT formulation cancels origin time uncertainty. \\
      \textbf{Magnitude Inversion} & Wood-Anderson ($r \le 30\,\unit{\kilo\meter}$) & Network-wide average & Mitigates bimodal underestimation on micro-earthquakes ($M_{\mathrm{L}} < 1.5$). \\
      \textbf{Hardware Backend} & Multi-GPU Blade ($4\times$ A100) & Multi-Core Workstation & Accelerates continuous multi-station inference to $>2{,}000\times$ real-time. \\
      \bottomrule
    \end{tabularx}
    \begin{tablenotes}[flushleft]
      \footnotesize
      \item Target \& Scope: Concrete operational specification for implementing AI seismic pipelines in regional European networks.
      \item What the reader should notice: The architecture couples regional deep-learning inference with robust non-linear location and velocity-model-aware association.
      \item Evidence Anchor: Synthesized from empirical benchmarks established in Chapter~\ref{chap:results}.
    \end{tablenotes}
  \end{threeparttable}
\end{table}

% =============================================================================
% SECTION 5: LIMITATIONS AND FUTURE RESEARCH HORIZONS
% =============================================================================
\section{Limitations and Future Research Horizons}
\label{sec:conclusion_future_work}

While this doctoral research establishes a verified baseline for AI-driven regional seismology, several scientific boundaries and emerging opportunities chart the roadmap for future development:

\begin{enumerate}
  \item \textbf{Scaling to the Pan-European AdriaArray Mega-Experiment}:
  This study evaluated a 93-day window across 266 stations. The ongoing **AdriaArray** initiative integrates over $1{,}500$ broadband and temporary stations spanning 15 European nations across the Adriatic-Alpine-Pannonian realm. Extending our CINECA Leonardo HPC workflow to the full multi-year AdriaArray dataset will enable the construction of the most complete, unified transboundary microseismic catalog in European history.
  \item \textbf{Integration of High-Resolution 3D Velocity Models and Adjoint Tomography}:
  All relocations in this dissertation utilized a 1D layered crustal velocity model. While optimal for baseline comparisons against historical bulletins, 1D models cannot account for the severe lateral velocity contrasts across the Alpine Moho descent (where Moho depth varies from $25\,\unit{\kilo\meter}$ beneath the Venetian plain to $> 50\,\unit{\kilo\meter}$ beneath the Tauern Window). Inverting travel times through high-resolution 3D regional tomographic models using NonLinLoc 3D or Eikonal wavefield solvers will sharpen hypocentral alignments and depth constraints along active blind thrust faults.
  \item \textbf{Real-Time Edge Streaming and IoT Integration}:
  Transitioning from batch HPC processing to sub-second streaming inference requires compiling deep neural networks (PhaseNet and EQTransformer) to lightweight edge-computing runtimes (such as NVIDIA TensorRT, ONNX, or low-power microcontrollers collocated directly with station digitizers). Edge picking will enable localized decentralized trigger verification and reduce continuous satellite telemetry bandwidth.
  \item \textbf{Automated First-Motion Polarities and Focal Mechanism Inversion}:
  While our pipeline successfully picked arrival times and local magnitudes, automated earthquake source characterization requires determining fault-plane kinematics. Future work should integrate neural first-motion polarity classifiers (e.g., classifying impulsive compressional vs.\ dilatational onsets) coupled with amplitude-ratio inversion (HASH algorithm) or full-waveform regional moment tensor inversion to automatically deliver focal mechanism solutions for events with $M_{\mathrm{L}} \ge 2.5$ within minutes of origin.
  \item \textbf{Waveform Cross-Correlation and Double-Difference Relocation}:
  The automated catalogs generated by PyOcto and NonLinLoc provide an ideal starting baseline for high-precision relative relocation. Integrating waveform cross-correlation delay-time measurements and double-difference relocation (\ac{HypoDD}) will collapse diffuse hypocentral clouds into razor-sharp fault planes, illuminating the fine-scale segmentation of active thrust ramps in the Friuli Prealps.
\end{enumerate}

% =============================================================================
% SECTION 6: CONCLUDING REMARKS
% =============================================================================
\section{Concluding Remarks}
\label{sec:conclusion_remarks}

Earthquake monitoring in seismically hazardous, densely populated regions such as Northeastern Italy demands both physical rigor and rapid operational responsiveness. This dissertation provides a reviewable implementation and evaluation framework for that work. A defensible operational conclusion requires a reproducibility bundle for every benchmark, independent adjudication of proposed events, and human approval of the relevant observatory workflow; none of these decisions is replaced by this draft.

